\section{Conclusion}


In this paper, we propose a unified dual-mask physical model for light field depth estimation to address the severe geometric degradation caused by non-Lambertian reflectance and complex textures. Moving beyond the conventional Epipolar Plane Image (EPI) linearity assumption, our framework explicitly models the physical degradation of angular signals to recover accurate depth maps in heterogeneous environments. By integrating reflectance priors with geometric constraints, the proposed method effectively resolves the depth ambiguities inherent in multi-view imaging systems deployed in unconstrained urban and indoor scenes.

The specific contributions and findings of this work are summarized as follows:
\textit{ \textbf{Geometry-based three-layer angular signal decomposition and non-Lambertian detection:} We decouple mixed angular signals into diffuse, specular, and scattering components. This mechanism accurately identifies non-Lambertian regions by analyzing geometric deviations, preventing the erroneous depth assignments typically caused by dynamic highlight shifts and volumetric scattering that violate standard EPI assumptions.
} \textbf{Unified dual-mask physical model:} We construct a dual-mask architecture that processes Lambertian and non-Lambertian regions simultaneously. By embedding physical reflectance constraints into the mask generation process, the model mitigates the slope ambiguity inherent in traditional EPI-based methods, particularly in complex textured scenes where photo-consistency is compromised.
\textbf{Multi-domain balanced training strategy:} We introduce an optimization scheme designed to counteract the extreme data scarcity in non-Lambertian scenarios. This strategy stabilizes gradient updates across heterogeneous material domains, achieving an overall MAE of 0.133 and a Mixed (Urban) MAE of 0.081, validating the system's robustness in diverse real-world environments.

The established physical modeling paradigm provides a solid foundation for dense geometry recovery in scenes with heterogeneous materials. Extending these reflectance-aware mechanisms to dynamic light field video systems will further advance real-time 3D scene understanding and immersive rendering applications.


The transition from frequency-domain analysis to explicit time-domain physical decomposition addresses a fundamental bottleneck in practical light field systems: the angular Nyquist limit. Traditional methods relying on 2D Discrete Fourier Transform (DFT) implicitly assume dense angular sampling to resolve the spectral signatures of non-Lambertian materials. However, in real-world multi-view imaging systems constrained to low angular resolutions (e.g., $9 \times 9$ or lower), sparse sampling induces severe spectral aliasing and leakage. This renders frequency-based Epipolar Plane Image (EPI) slope extraction fundamentally unreliable. By formulating the three-layer angular signal decomposition as $I(u,v) = DC + a\cdot u + b\cdot v + \varepsilon(u,v)$, our approach systematically bypasses this Nyquist constraint. The residual term $\varepsilon(u,v)$ explicitly captures high-frequency material variations directly in the spatial-angular domain. This demonstrates that explicit geometric decoupling is inherently more robust and computationally efficient than implicit spectral filtering when the angular bandwidth of the capture system is strictly limited.

Furthermore, the integration of physics-based geometric features with a Random Forest (RF) classifier highlights a critical system-level trade-off between representational capacity and few-shot generalization. While end-to-end Convolutional Neural Networks (CNNs) excel at implicit feature learning, they notoriously overfit when non-Lambertian training data is scarce, which is typical for unconstrained urban datasets. Our empirical findings—where the angular gradient feature yields a Cohen's $d$ effect size of 0.983 and the RF classifier achieves an AUC of 0.962—validate that handcrafted physical priors provide superior out-of-distribution robustness. This interpretability is crucial for downstream 3D video and autonomous navigation systems. It allows the GeometricDualMask network to route features adaptively based on verifiable physical properties, such as parallax magnitude and material ratio, rather than relying on opaque latent embeddings. Consequently, this ensures stable and reliable depth regression across heterogeneous environments without requiring massive, fully annotated non-Lambertian light field datasets.

Despite these structural advantages, the current framework exhibits limitations in scenarios involving severe multi-view occlusions or extreme specularities where the foundational angular cone constraint completely collapses. When the non-Lambertian residual $\varepsilon(u,v)$ overwhelmingly dominates the linear parallax terms, the extracted geometric features may degrade, potentially leading to suboptimal mask generation. Additionally, while the domain-balanced sampling strategy effectively mitigates the long-tail distribution of complex materials, it introduces non-trivial computational overhead during the training phase. Future work will investigate integrating this dual-mask physical prior into modern differentiable rendering pipelines, such as 3D Gaussian Splatting or Neural Radiance Fields, to jointly optimize dense geometry and spatially-varying BRDFs. Furthermore, extending the angular signal decomposition to accommodate transient imaging or time-of-flight light fields could ultimately bridge the remaining gap between physical light transport modeling and robust, real-time 3D scene understanding.

Despite the demonstrated efficacy of the proposed dual-mask physical model in resolving depth ambiguities, several limitations remain. First, the three-layer angular signal decomposition, formulated as $I(u,v) = DC + a\cdot u + b\cdot v + \epsilon(u,v)$, inherently relies on a local linearity assumption. While highly effective for isotropic specularities and standard volumetric scattering, this formulation struggles with highly anisotropic Bidirectional Reflectance Distribution Functions (BRDFs), such as brushed metals or complex woven fabrics. In such cases, the residual term $\epsilon$ may fail to fully decouple intricate directional reflectance, potentially leading to mask misclassification. Second, although our geometric feature extraction is explicitly optimized for low angular resolution light fields (e.g., $9 \times 9$), the computation of angular gradients and correlation coherence becomes mathematically ill-posed in extremely sparse setups (e.g., $2 \times 2$ or $3 \times 3$ multi-view configurations). Moreover, severe occlusions in highly cluttered urban scenes can disrupt Epipolar Plane Image slope continuity, causing geometric features to conflate occlusion boundaries with non-Lambertian transitions. Finally, the current framework processes each light field frame independently, lacking temporal regularization, which may introduce noticeable flickering artifacts in non-Lambertian regions when deployed in dynamic video light field systems.

To address these challenges, future work will explore integrating our physical priors with emerging neural rendering paradigms, such as 3D Gaussian Splatting or Neural Radiance Fields. By embedding the dual-mask routing mechanism into the view-dependent color networks of implicit representations, we aim to jointly optimize depth estimation and novel view synthesis for scenes featuring extreme anisotropic materials. Additionally, extending the framework to dynamic light field videos is a critical next step for immersive virtual reality applications. This will involve incorporating spatio-temporal consistency constraints and optical flow-guided mask propagation to suppress temporal jitter in specular and translucent regions. Another promising direction is transitioning from supervised mask generation to self-supervised learning, leveraging physical consistency losses to eliminate the reliance on scarce ground-truth non-Lambertian annotations. Furthermore, we plan to investigate hardware-algorithm co-design strategies. By collaborating with computational photography techniques, such as coded apertures or metasurface lenses, we can optically encode non-Lambertian signatures directly at the sensor level. This would inherently enhance the signal-to-noise ratio of the residual term $\epsilon$, thereby reducing the computational overhead of the backend classifier and enabling real-time, high-fidelity depth recovery on resource-constrained edge devices.